\histitem{Compact operators}

Riesz considerably amplifies these remarks less than three years later, in his memoir on functional equations [60] which develops, with a clarity undimmed after more than a century, the essentials of the spectral theory of compact operators on Banach spaces. The stated aim of this text is to reformulate Fredholm theory for equation (1) on a compact interval I, by considering the space \(E = \mathcal{C}(I)\) equipped with the topology of uniform convergence. But as we have already indicated in an earlier note (ÉHM, p.~268), Riesz is clearly aware that his results apply to any Banach space — a notion that would only be introduced in the following decade.

Inspired by Fréchet's ideas on general topology, Riesz now considers the endomorphisms of E that map every bounded sequence to a relatively compact sequence. Like Hilbert, he still calls them "completely continuous" and notes that this new notion is equivalent, in the case of \(\mathrm{E} = \ell_{\mathbf{C}}^{2}(\mathbf{N})\) , to the notion of complete continuity that he had used in his 1913 book \(^{(4)}\) .

Riesz studies endomorphisms of the form \(B = 1_{E} - A\) where A is a completely continuous endomorphism of E, and derives all their spectral properties from the fact that a locally compact normed vector space must be finite-dimensional — a fact apparently discovered for the occasion. He proves easily that the kernel of B is finite-dimensional and that its image is closed and of finite codimension. As E. Weyr [87] had done it in finite dimension, he then considers the iterates \(B^{k}\) . He shows that the sequence of their kernels and that of their images are stationary, then introduces what we today call the nilspace and conilspace of B, efficiently showing that E is their topological direct sum. Applying this observation to endomorphisms \(B_{\lambda} = 1_{E} - \lambda A\) for \(\lambda \in C\) , he deduces without difficulty the spectral properties of completely continuous endomorphisms of a Banach space, and this in a more or less definitive form.

Among the main results on the spectrum of these operators, only those that would require a more mature view of the notion of function space, particularly of duality, seem to escape Riesz. For example, in the late 1920s, T. Hildebrandt [33] and J. Schauder [63] will show (within the now well-established framework of Banach spaces) that the adjoint of a completely continuous endomorphism is completely continuous.

Moreover, the notion of Riesz complete continuity still depends on the use of sequences; Banach [4] will free himself from it in his 1932 book to consider mappings that transform every bounded set into a relatively compact set. The term “compact” linear mapping would only gradually establish itself in the second half of the 20th century, presumably following E. Hille's book [34] on semigroups of operators.

Several questions raised in the 1930s would remain open for a long time. In this direction, we note the resolution [18] in 1973 of the "approximation problem" made famous by Banach and S. Mazur in 1932 and 1938 (see Remark 6 of III, p.~16 and Exercise 25 of III, p.~112).

At the same time, Lomonosov [40] proves the existence of nontrivial closed invariant subspaces for an endomorphism commuting with a nonzero compact endomorphism in a locally convex space (cor. 2 of III, p.~13). This result was one of the most significant advances on the problem of the existence of such a subspace for an arbitrary continuous endomorphism (the “invariant subspace problem”). This problem is still open at present in the case of continuous endomorphisms of separable Hilbert spaces; P. Enflo [19] constructed in 1987 the first example of an endomorphism of a Banach space admitting no nontrivial closed invariant subspace.
% semantic-fragments: legacy-input-seam
\relax{}
